\section{Proposed Architecture}\label{seec: proposed architecture}

\subsection{Overview}\label{subsec: overview}

While approaches in the literature propose to learn a monolithic approximator $\predMono = \nnMono(\vector{\rotAngle{}}):\realSet{\nJoints}\rightarrow\realSet{\nDims\times\nJoints}$, our main idea is to learn one approximator for each joint $\joint{i}$, which encodes $\matrix{T}_{i-1}^{i}$, and thus learning $\nJoints$ predictors $\predLink_i = \nnLink_i(\vector{\rotAngle{i}}):\realSet{1}\rightarrow\realSet{\nDims\times(\nDims+1)}~|~i\in\{1,\hdots\nJoints\}$, and joining all $\nnLink_i$ leads to the proposed model $\nnLink$.


\begin{figure}
	\centering
	\includegraphics[width=\linewidth]{linked_network_overview.pdf}
	\caption{Proposed architecture. For each joint $i\in\{1,\hdots,j\}$, a network $\nnLink_i$ learns the \gls{dh} transformation from joint $i-1$ to $i$ $\matrix{R}_{i-1}^{i}$.
	%
	$\matrix{R}_{i-1}^{i}$ is multiplied by the cumulative transformation of all previous joints, resulting in $\matrix{R}_{0}^{i}$, from which the Cartesian position of joint $i$ $\prediction{\vector{x}}_i$ is obtained directly.}
	\label{fig:linked network overview}
\end{figure}

\subsection{Loss Function}\label{subsec: loss}

For training $\nn$, we follow the approaches in the literature \rev{citation needed} and use the average Euclidean distance between the joints and their respective prediction, as described in Equation \ref{eq: loss}.
%
Note that some approaches in the literature consider only the end-effector position and its prediction, which is reasonable since in a real robot, the position of individual joints is hard to obtain, while the end-effector's position can be obtained with off-the-shelf trackers and cameras \rev{citation}.
%
In contrast, since we perform the training in simulation, we have the flexibility of considering all joints, leading to a richer loss signal, improving the training convergence.

\begin{equation}\label{eq: loss}
	\loss = \frac{1}{\nJoints}\sum_{k=1}^{\nJoints}||\prediction{x}_k-\vector{x}_k||
\end{equation}

%Similarly, for each $\nnLink_i$, we average the Euclidean distance between for all joints up to $\joint{i}$, according to Equation \ref{eq: loss link cart}.
%\begin{equation}\label{eq: loss link cart}
%	\lossLinkCart_i = \frac{1}{i}\sum_{k=1}^{i}||\predMono_k-\vector{x}_k|| ~|~ i\in\{1,\hdots,j\}
%\end{equation}
%Furthermore, we explore two properties of rotation matrices by adding two extra loss terms for $\nnLink_i$.
%
%The first property is that $\det(\matrix{R}) = 1$, which is captured by Equation \ref{eq: loss link det}.
%
%The second property is that $\matrix{R}^{T} = \matrix{R}^{-1}$, which is capture by Equation \ref{eq: loss link inv}.
%\begin{equation}\label{eq: loss link det}
%	\lossLinkDet_i = ||\det(\prediction{\matrix{R}}_i) - 1|| ~|~ i\in\{1,\hdots,j\}
%\end{equation}
%
%\begin{equation}\label{eq: loss link inv}
%	\lossLinkInv_i = ||\prediction{\matrix{R}}_i\prediction{\matrix{R}}_i^{T} - \matrix{I}|| ~|~ i\in\{1,\hdots,j\}
%\end{equation}
%The final loss for training $\nnLink_i$ is given by Equation \ref{eq: loss link}.
%
%The reason for incorporating $\lossLinkInv_i$ and $\lossLinkDet_i$ in $\lossLink_i$ is to motivate $\nnLink_i$ to converge to valid \gls{dh} matrices, indirectly pushing $\prediction{\matrix{D}}_i$ to converge to the proper Cartesian position of $\joint{i}$. 
%\begin{equation}\label{eq: loss link}
%	\lossLink_i = \lossLinkCart_i+\lossLinkInv_i+\lossLinkDet_i
%\end{equation}

\subsection{Training Procedure}\label{subsec: training procedure}

Since in this stage we use simulation, both $\nnMono$ and $\nnLink$ are trained in a straight-forward 
fashion according to Algorithm \ref{alg: train nn}.
%
For each epoch of training, $nSamples$ random poses are sampled in joint space (line $3$), and the respective Cartesian position of each joint is obtained directly from the simulator (line $4$) and added to a memory buffer (line $5$).
%
The memory size is limited, and upon additions further than its limit, it drops a percentage of its oldest data (we use a $10\%$ data drop in all experiments) for reducing data allocation overheads.
%
Next, $nIterations$ of training is done with the data already within the memory, which are done by randomly sampling $batchSize$ samples from the memory (line $8$), computing the average loss (lines $9$ to $13$), which is then back-propagated through the network (line $14$).
 
\begin{algorithm}
	\caption{Training procedure for $\nn$.}\label{alg: train nn}
	\For{$e \in \{1,\hdots, epochs\}$}{
		\For{$r \in \{1,\hdots, nSamples\}$}{
			$q\leftarrow$ random joint positions\;
			$x\leftarrow$ Cartesian position of all joints given $q$\;
			add $(q, x)$ to $memory$\;
		}
		
		\For{$t\in\{1,\hdots,nIterations\}$}{
			$(qs, xs)\leftarrow$ get $batchSize$ random samples from $memory$\; 
			$meanLoss\leftarrow0$\;
			\For{$(q, x) \in (qs, xs)$}{
				$meanLoss\leftarrow acc+\loss$\;
			}
			$meanLoss\leftarrow acc/batchSize$\;
			backpropagate $meanLoss$ through $\nn$\; 
		}
	}
\end{algorithm}

It is worthy to notice that in our tests, we observed that randomly sampling poses leads to a much higher efficiency than performing random walks, which is to be expected due to the better state space coverage. 
%
However, when considering a real robot in the \gls{rl} loop, it is unpractical to sample the state space this way, since the robot will be naturally performing a random walk.
%
In section \ref{sec: rl case of study}, we present a \gls{rl} case of study highlighting the proposed architecture impacts in this scope.